\subsection{Fringe Movement to Mirror Displacement and Velocity}
\label{sec:fringe_to_displacement_velocity}

A detector only delivers intensity values; to obtain a physical quantity such as a displacement or a velocity, those values must be linked to the quantity of interest through a model. In interferometry this link is the fringe count: since Michelson used it to express the meter in wavelengths of light, the passage of fringes has been the way of converting a change of optical path into a length. Establishing this relation is what turns the interferometer and its detectors from an instrument that produces a pattern into an instrument that measures.

For this work the theoretical background is relevant for several reasons. It fixes the displacement associated with each fringe, and     therefore the resolution that fringe counting can reach (\cref{eq:displacement_m2}). It relates the velocity of the mirror to the rate at which fringes cross the detector (\cref{eq:velocity_fringe_rate}), which is the quantity that limits how fast the mirror can move before the camera starts to lose fringes (\cref{sec:camera_velocity}). It also shows that a single intensity value cannot tell in which direction the mirror moves, since the sign of the motion is encoded in the direction of the fringe progression, which motivates the use of the camera and its spatial information.

The following sections derive the relation between fringe movement, mirror displacement and velocity, and then present the method used to apply it to the frames of the camera by tracking the irradiance of a region of interest (\cref{sec:roi_fringe_velocity}), together with the graphical interface that implements it (\cref{sec:camera_gui}).

Let $L_2(t)$ denote the optical path length from $\text{BS}$ to the movable mirror $M2$ at time $t$, and let $L_1$ denote the constant optical path length along the fixed path to $M1$. Because each beam traverses the distance between $\text{BS}$ and its corresponding mirror twice, the round-trip optical path difference in a medium of refractive index $n$ (with $n \approx 1$ in air) is
\begin{equation}
    \Delta(t) = 2n\bigl[L_2(t) - L_1\bigr] + \Delta_0
    \label{eq:opd_m2}
\end{equation}
where $\Delta_0$ accounts for any constant phase-shift offset (such as reflection phase differences at $\text{BS}$ and the mirrors).

\subsubsection{Mirror Displacement from Fringe Shift}
Suppose the shift of $\Delta m = 1$ corresponds to a mirror displacement of $\Delta d_{M2} = \lambda/2$ along the axis of beam $(2)$. Tracking continuous phase $\delta(t)$ rather than integer fringe counts alone also permits sub-fringe displacement measurement.

\subsubsection{Mirror Velocity from Fringe Shift}
Differentiating \cref{eq:displacement_m2} with respect to time yields the velocity $v_{M2}(t)$ of the movable mirror $M2$:
\begin{equation}
    v_{M2}(t) = \frac{\lambda}{2}\,\frac{dm(t)}{dt}
    \label{eq:velocity_m2}
\end{equation}

In terms of the fringe-crossing rate (temporal fringe frequency) $f_{\text{fringe}}(t) = \frac{dm(t)}{dt}$ observed at a fixed detector or camera ROI along output beam $(4)$, the mirror speed is
\begin{equation}
    v_{M2}(t) = \frac{\lambda}{2n}\,f_{\text{fringe}}(t)
    \label{eq:velocity_fringe_rate}
\end{equation}


Alternatively, if a spatial fringe feature shifts across a coordinate axis $x$ perpendicular to the fringe normal with spatial fringe period $\Lambda_x$ and fringe-shift speed $v_{\text{fringe}}(t) = \frac{dx(t)}{dt}$, then $f_{\text{fringe}}(t) = \frac{v_{\text{fringe}}(t)}{\Lambda_x}$, giving
\begin{equation}
    v_{M2}(t) = \frac{\lambda}{2n}\,\frac{v_{\text{fringe}}(t)}{\Lambda_x}
    \label{eq:velocity_spatial_shift}
\end{equation}


The sign of $\Delta d_{M2}$ and $v_{M2}(t)$ is determined by the direction of fringe progression, which corresponds to $M2$ moving toward or away from $\text{BS}$ along beam $(2)$. The movable mirror $M2$ translates along the optical axis of beam $(2)$ by a displacement $\Delta d_{M2} = d_{M2}(t_2) - d_{M2}(t_1)$ between two instants $t_1$ and $t_2$. This changes the round-trip optical path length by $2\,\Delta d_{M2}$, producing a corresponding change in the interference phase difference $\delta = \frac{2\pi}{\lambda}\Delta$:
\begin{equation}
    \Delta\delta = \frac{4\pi}{\lambda}\,\Delta d_{M2}
    \label{eq:phase_shift_m2}
\end{equation}

Where $\lambda$ is the wavelength of the laser source.


Because each full fringe cycle (from one bright maximum to the next, or one dark minimum to the next) corresponds to a phase change of $\Delta\delta = 2\pi\,\Delta m$, where $\Delta m$ is the net number of fringes passing a fixed observation point or region of interest (ROI), substituting into \cref{eq:phase_shift_m2} gives the axial displacement of mirror $M2$:
\begin{equation}
    \Delta d_{M2} = \frac{\lambda}{2}\,\Delta m
    \label{eq:displacement_m2}
\end{equation}


Each fringe shift of $\Delta m = 1$ corresponds to a mirror displacement of $\Delta d_{M2} = \lambda/2$ along the axis of beam $(2)$. Tracking continuous phase $\delta(t)$ rather than integer fringe counts alone also permits sub-fringe displacement measurement.